\section{Analisis Keamanan}
\label{sec:keamanan-dh}

Keamanan Diffie--Hellman bersandar sepenuhnya pada kesulitan menghitung
logaritma diskret dalam grup perkalian bilangan bulat modulo $p$. Bagian ini
membahas seberapa kuat sandaran itu, dan kemudian memperlihatkan bahwa kekuatan
tersebut tidak menolong sama sekali terhadap serangan yang bersifat aktif.

\subsection{Masalah Logaritma Diskret}
\label{subsec:dlog-dh}

Diberikan bilangan prima $p$, akar primitif $g$, dan sebuah nilai $A$ dengan
$1 \le A \le p-1$, masalah logaritma diskret menuntut pencarian bilangan bulat
$a$ sehingga

\begin{equation}
g^a \equiv A \pmod p .
\label{eq:dlp}
\end{equation}

Berbeda dari logaritma biasa, tidak ada rumus tertutup untuk
Persamaan~\eqref{eq:dlp}. Nilai $a$ hanya dapat dicari dengan menjelajahi
ruang kemungkinan, sehingga seluruh pertaruhannya adalah \emph{kecepatan}
penjelajahan tersebut.

\paragraph{Pencarian menyeluruh.} Cara paling lugu adalah mencoba
$a = 1, 2, 3, \dots$ sampai Persamaan~\eqref{eq:dlp} terpenuhi. Cara itu
memerlukan paling banyak $p-1$ percobaan, sehingga untuk modulus 2048 bit
diperlukan sekitar $2^{2048}$ percobaan --- angka yang tidak ada artinya secara
fisik.

\paragraph{Langkah raksasa-bayi (\textit{baby-step giant-step}).} Penjelajahan
dapat dipercepat secara mendasar dengan menyimpan hasil-hasil antara. Tuliskan
$a = i \cdot m + j$ dengan $m = \lceil \sqrt{p-1} \rceil$, sehingga
Persamaan~\eqref{eq:dlp} menjadi

\[ g^{\,i m + j} \equiv A \pmod p
   \quad\Longleftrightarrow\quad
   g^{\,j} \equiv A \cdot \left(g^{-m}\right)^{i} \pmod p . \]

Langkah pertama menyimpan $g^j$ untuk seluruh $0 \le j < m$ ke dalam tabel;
langkah kedua mengalikan $A$ dengan $g^{-m}$ berulang kali sampai hasilnya
ditemukan di dalam tabel. Karena setiap langkah berukuran $\sqrt{p}$, waktu
maupun memori yang diperlukan hanya sekitar $2\sqrt{p}$.

Pada parameter kecil $p = 353$ dan $g = 3$, akar kuadrat dari 352 adalah sekitar
18,8 sehingga $m = 19$. Algoritma itu memerlukan 19 langkah untuk menemukan
$a = 5 \cdot 19 + 2 = 97$, padahal pencarian menyeluruh memerlukan 97 langkah.
Perbaikan itu terasa kecil pada modulus sepanjang 9 bit, tetapi pada modulus
2048 bit ia mengubah $2^{2048}$ menjadi sekitar $2^{1024}$ --- perbaikan yang
sangat besar, tetapi angka hasilnya masih mustahil. Perhatikan bahwa percepatan
ini juga menuntut memori sebesar $\sqrt{p}$, dan pada modulus besar kebutuhan
memori itulah yang lebih dahulu menjadi penghalang.

\paragraph{Metode yang lebih baik.} Untuk medan prima, algoritma terbaik yang
dikenal hari ini bekerja dalam waktu \textit{subeksponensial} terhadap panjang
modulus, dengan bentuk yang menyerupai hasil untuk pemfaktoran pada
Bab~\ref{chap:rsa}:

\begin{equation}
L_p\!\left[\tfrac{1}{3},\, (64/9)^{1/3}\right]
= \exp\!\left( (64/9)^{1/3} (\ln p)^{1/3} (\ln \ln p)^{2/3} \right) .
\label{eq:dlp-index-calculus}
\end{equation}

Karena bentuk fungsinya sama, panjang modulus yang memberi tingkat keamanan
setara ternyata juga berdekatan, sebagaimana diperlihatkan pada
Tabel~\ref{tab:ukuran-dh}. Angka-angka itu bukan hasil hitungan langsung,
melainkan rekomendasi baku yang diturunkan dari perkiraan biaya komputasi
terkini; sifatnya perkiraan, dan sebagaimana ditunjukkan pengalaman pemfaktoran
RSA, perkiraan semacam itu cenderung terlalu optimistis.

\begin{table}[htbp]
    \centering
    \begin{tabularx}{\linewidth}{@{}>{\raggedright\arraybackslash}p{2.0cm}
                                   >{\raggedright\arraybackslash}p{3.4cm}
                                   >{\raggedright\arraybackslash}X@{}}
    \toprule
    \textbf{Modulus $p$} & \textbf{Perkiraan keamanan} & \textbf{Keterangan} \\
    \midrule
    512 bit  & kurang dari 64 bit & pernah dipecahkan untuk parameter ekspor TLS \\
    1024 bit & sekitar 80 bit     & dianggap dalam jangkauan penyerang yang sangat kuat \\
    2048 bit & sekitar 112 bit    & panjang minimum yang diterima hari ini \\
    3072 bit & sekitar 128 bit    & untuk perlindungan jangka panjang \\
    7680 bit & sekitar 192 bit    & setara AES-192 \\
    15360 bit & sekitar 256 bit   & setara AES-256 \\
    \bottomrule
    \end{tabularx}
    \caption{Panjang modulus Diffie--Hellman dan perkiraan tingkat keamanan yang
    diberikannya. Perhatikan bahwa modulusnya jauh lebih panjang daripada
    modulus RSA yang memberi keamanan setara}
    \label{tab:ukuran-dh}
\end{table}

\subsection{Serangan Man-in-the-Middle}
\label{subsec:mitm-dh}

Kesulitan logaritma diskret melindungi Diffie--Hellman dari penyadapan, tetapi
tidak melindunginya dari apa pun yang lebih aktif daripada penyadapan. Serangan
yang paling penting adalah \textbf{man-in-the-middle} (MITM), yang pada dasarnya
memanfaatkan satu kelemahan mendasar: protokol ini tidak pernah memeriksa
\emph{siapa} yang berada di ujung lain percakapan.

Misalkan Mallory mampu mencegat dan mengganti pesan pada saluran antara Alice
dan Bob. Ia memilih kunci privatnya sendiri, $m$, dan menjalankan protokol itu
dua kali: sekali berpura-pura menjadi Bob di hadapan Alice, dan sekali
berpura-pura menjadi Alice di hadapan Bob. Urutan pertukarannya diperlihatkan
pada Tabel~\ref{tab:alur-mitm}.

\begin{table}[htbp]
    \centering
    \begin{tabular}{@{}l l l@{}}
    \toprule
    \textbf{Langkah} & \textbf{Alice} & \textbf{Bob} \\
    \midrule
    1 & mengirim $g^a \bmod p$ & (dicegat Mallory) \\
    2 & menerima $g^m$, dikira dari Bob & mengirim $g^b \bmod p$ \\
    3 & (dicegat Mallory) & menerima $g^m$, dikira dari Alice \\
    4 & menghitung $K_1 = (g^m)^a$ & menghitung $K_2 = (g^m)^b$ \\
    \bottomrule
    \end{tabular}
    \caption{Alur pertukaran kunci yang diserang Mallory. Nilai $g^a$ dan $g^b$
    yang asli tidak pernah sampai ke tujuannya; keduanya diganti dengan $g^m$}
    \label{tab:alur-mitm}
\end{table}

Kunci yang dipegang masing-masing pihak berbeda, dan Mallory mengetahui
keduanya:

\begin{align}
K_1 &= g^{am} \bmod p = (g^a)^m \bmod p = (g^m)^a \bmod p , \label{eq:mitm-k1} \\
K_2 &= g^{bm} \bmod p = (g^b)^m \bmod p = (g^m)^b \bmod p . \label{eq:mitm-k2}
\end{align}

Pada Persamaan~\eqref{eq:mitm-k1}, suku $(g^a)^m$ diperoleh Mallory dari kunci
publik Alice, sedangkan suku $(g^m)^a$ adalah nilai yang dihitung Alice sendiri.
Keduanya sama persis, dan itulah yang membuat serangan ini berhasil: Mallory
mendapatkan $K_1$ tanpa perlu mengetahui $a$, sebab ia memangkatkan $g^a$ dengan
kunci privatnya sendiri. Alasan yang sama berlaku untuk
Persamaan~\eqref{eq:mitm-k2}. Susunan ketiga pihak itu diperlihatkan pada
Gambar~\ref{fig:mitm-dh}.

% Serangan man-in-the-middle pada pertukaran kunci Diffie-Hellman.
% Dirujuk dari section-03.tex dengan % Serangan man-in-the-middle pada pertukaran kunci Diffie-Hellman.
% Dirujuk dari section-03.tex dengan % Serangan man-in-the-middle pada pertukaran kunci Diffie-Hellman.
% Dirujuk dari section-03.tex dengan \input{figures/mitm-dh}.
% Mallory mengganti g^a dan g^b dengan g^m, sehingga ia memegang kedua kunci.
\begin{figure}[htbp]
    \centering
    \resizebox{0.98\linewidth}{!}{%
    \begin{tikzpicture}[
        kotak/.style={draw, rounded corners=1pt, minimum height=0.9cm,
                      text width=2.7cm, align=center, font=\small, inner sep=3pt},
        panah/.style={-{Latex[length=2.2mm]}, thick},
        palsu/.style={-{Latex[length=2.2mm]}, thick, dashed},
        catatan/.style={font=\footnotesize, align=left, text width=13.6cm},
    ]
        % --- tiga pihak ---
        \node[kotak, fill=blue!6]   (a) at (0,0)    {Alice};
        \node[kotak, fill=gray!5]   (m) at (6.4,0)  {Mallory};
        \node[kotak, fill=green!10] (b) at (12.8,0) {Bob};

        % --- kunci yang dipegang masing-masing ---
        \node[font=\footnotesize, align=center] at (0,-1.2)
            {memegang $a$\\ memakai $K_1 = g^{am}$};
        \node[font=\footnotesize, align=center] at (6.4,-1.2)
            {memegang $m$\\ mengetahui $K_1$ dan $K_2$};
        \node[font=\footnotesize, align=center] at (12.8,-1.2)
            {memegang $b$\\ memakai $K_2 = g^{bm}$};

        % --- pengiriman yang dicegat ---
        \draw[panah] (a.east) -- node[above, font=\scriptsize] {$g^{a} \bmod p$} (2.6,0);
        \node[font=\scriptsize, align=center] at (4.5,0.42) {diganti};
        \draw[palsu] (m.west) -- node[below, font=\scriptsize] {$g^{m} \bmod p$} (2.6,0);

        \draw[palsu] (m.east) -- node[above, font=\scriptsize] {$g^{m} \bmod p$} (10.2,0);
        \node[font=\scriptsize, align=center] at (11.5,0.42) {seolah dari Alice};
        \draw[panah] (b.west) -- node[below, font=\scriptsize] {$g^{b} \bmod p$} (10.2,0);

        \node[catatan, anchor=north] at (6.4,-2.2)
            {Alice dan Bob masing-masing merasa berbicara dengan pihak yang benar,
             sebab kedua perhitungan berjalan mulus. Mallory memakai $K_1$ untuk
             membaca arah Alice dan $K_2$ untuk membaca arah Bob};
    \end{tikzpicture}}
    \caption{Serangan man-in-the-middle: $g^m$ menggantikan $g^a$ dan $g^b$,
    sehingga Mallory menyepakati kunci yang berbeda dengan masing-masing pihak}
    \label{fig:mitm-dh}
    \sumber{Diolah dari Munir, \textit{IF4020 Kriptografi}, ITB.}
\end{figure}
.
% Mallory mengganti g^a dan g^b dengan g^m, sehingga ia memegang kedua kunci.
\begin{figure}[htbp]
    \centering
    \resizebox{0.98\linewidth}{!}{%
    \begin{tikzpicture}[
        kotak/.style={draw, rounded corners=1pt, minimum height=0.9cm,
                      text width=2.7cm, align=center, font=\small, inner sep=3pt},
        panah/.style={-{Latex[length=2.2mm]}, thick},
        palsu/.style={-{Latex[length=2.2mm]}, thick, dashed},
        catatan/.style={font=\footnotesize, align=left, text width=13.6cm},
    ]
        % --- tiga pihak ---
        \node[kotak, fill=blue!6]   (a) at (0,0)    {Alice};
        \node[kotak, fill=gray!5]   (m) at (6.4,0)  {Mallory};
        \node[kotak, fill=green!10] (b) at (12.8,0) {Bob};

        % --- kunci yang dipegang masing-masing ---
        \node[font=\footnotesize, align=center] at (0,-1.2)
            {memegang $a$\\ memakai $K_1 = g^{am}$};
        \node[font=\footnotesize, align=center] at (6.4,-1.2)
            {memegang $m$\\ mengetahui $K_1$ dan $K_2$};
        \node[font=\footnotesize, align=center] at (12.8,-1.2)
            {memegang $b$\\ memakai $K_2 = g^{bm}$};

        % --- pengiriman yang dicegat ---
        \draw[panah] (a.east) -- node[above, font=\scriptsize] {$g^{a} \bmod p$} (2.6,0);
        \node[font=\scriptsize, align=center] at (4.5,0.42) {diganti};
        \draw[palsu] (m.west) -- node[below, font=\scriptsize] {$g^{m} \bmod p$} (2.6,0);

        \draw[palsu] (m.east) -- node[above, font=\scriptsize] {$g^{m} \bmod p$} (10.2,0);
        \node[font=\scriptsize, align=center] at (11.5,0.42) {seolah dari Alice};
        \draw[panah] (b.west) -- node[below, font=\scriptsize] {$g^{b} \bmod p$} (10.2,0);

        \node[catatan, anchor=north] at (6.4,-2.2)
            {Alice dan Bob masing-masing merasa berbicara dengan pihak yang benar,
             sebab kedua perhitungan berjalan mulus. Mallory memakai $K_1$ untuk
             membaca arah Alice dan $K_2$ untuk membaca arah Bob};
    \end{tikzpicture}}
    \caption{Serangan man-in-the-middle: $g^m$ menggantikan $g^a$ dan $g^b$,
    sehingga Mallory menyepakati kunci yang berbeda dengan masing-masing pihak}
    \label{fig:mitm-dh}
    \sumber{Diolah dari Munir, \textit{IF4020 Kriptografi}, ITB.}
\end{figure}
.
% Mallory mengganti g^a dan g^b dengan g^m, sehingga ia memegang kedua kunci.
\begin{figure}[htbp]
    \centering
    \resizebox{0.98\linewidth}{!}{%
    \begin{tikzpicture}[
        kotak/.style={draw, rounded corners=1pt, minimum height=0.9cm,
                      text width=2.7cm, align=center, font=\small, inner sep=3pt},
        panah/.style={-{Latex[length=2.2mm]}, thick},
        palsu/.style={-{Latex[length=2.2mm]}, thick, dashed},
        catatan/.style={font=\footnotesize, align=left, text width=13.6cm},
    ]
        % --- tiga pihak ---
        \node[kotak, fill=blue!6]   (a) at (0,0)    {Alice};
        \node[kotak, fill=gray!5]   (m) at (6.4,0)  {Mallory};
        \node[kotak, fill=green!10] (b) at (12.8,0) {Bob};

        % --- kunci yang dipegang masing-masing ---
        \node[font=\footnotesize, align=center] at (0,-1.2)
            {memegang $a$\\ memakai $K_1 = g^{am}$};
        \node[font=\footnotesize, align=center] at (6.4,-1.2)
            {memegang $m$\\ mengetahui $K_1$ dan $K_2$};
        \node[font=\footnotesize, align=center] at (12.8,-1.2)
            {memegang $b$\\ memakai $K_2 = g^{bm}$};

        % --- pengiriman yang dicegat ---
        \draw[panah] (a.east) -- node[above, font=\scriptsize] {$g^{a} \bmod p$} (2.6,0);
        \node[font=\scriptsize, align=center] at (4.5,0.42) {diganti};
        \draw[palsu] (m.west) -- node[below, font=\scriptsize] {$g^{m} \bmod p$} (2.6,0);

        \draw[palsu] (m.east) -- node[above, font=\scriptsize] {$g^{m} \bmod p$} (10.2,0);
        \node[font=\scriptsize, align=center] at (11.5,0.42) {seolah dari Alice};
        \draw[panah] (b.west) -- node[below, font=\scriptsize] {$g^{b} \bmod p$} (10.2,0);

        \node[catatan, anchor=north] at (6.4,-2.2)
            {Alice dan Bob masing-masing merasa berbicara dengan pihak yang benar,
             sebab kedua perhitungan berjalan mulus. Mallory memakai $K_1$ untuk
             membaca arah Alice dan $K_2$ untuk membaca arah Bob};
    \end{tikzpicture}}
    \caption{Serangan man-in-the-middle: $g^m$ menggantikan $g^a$ dan $g^b$,
    sehingga Mallory menyepakati kunci yang berbeda dengan masing-masing pihak}
    \label{fig:mitm-dh}
    \sumber{Diolah dari Munir, \textit{IF4020 Kriptografi}, ITB.}
\end{figure}



Setelah itu Mallory berada dalam kedudukan yang sangat menguntungkan. Pesan
terenkripsi yang dikirim Alice dengan $K_1$ dapat ia dekripsi, ia baca, ia ubah
bila perlu, lalu ia enkripsi kembali dengan $K_2$ dan teruskan kepada Bob.
Sebaliknya, pesan dari Bob dapat ia buka dengan $K_2$ dan teruskan dengan $K_1$.
Alice dan Bob masing-masing merasa berkomunikasi langsung dengan pihak yang
benar, karena dari sudut pandang keduanya seluruh perhitungan berjalan mulus dan
menghasilkan kunci yang konsisten.

Pada parameter yang dipakai di seluruh bab ini, yaitu $p = 353$ dan $g = 3$
dengan $a = 97$ dan $b = 233$, seorang Mallory yang memilih $m = 101$ akan
memperoleh $g^m \bmod 353 = 63$. Alice lalu menghitung
$K_1 = 63^{97} \bmod 353 = 257$, sedangkan Bob menghitung
$K_2 = 63^{233} \bmod 353 = 249$. Mallory dapat memperoleh kedua nilai itu
dengan menghitung $40^{101} \bmod 353$ dan $248^{101} \bmod 353$, memakai
kunci publik yang ia sadap. Perhatikan bahwa $K_1 \neq K_2$; keduanya memang
harus berbeda, sebab bila sama maka Mallory tidak dapat memisahkan kedua arah
percakapan.

\subsection{Mengapa Serangan Itu Berhasil dan Bagaimana Menangkalnya}
\label{subsec:penangkalan-mitm}

Akar penyebab serangan MITM bukanlah kelemahan matematis, melainkan kelemahan
protokol: tidak ada satu pun langkah pada Subbagian~\ref{sec:protokol-dh} yang
memeriksa identitas lawan bicara. Protokol itu menjamin bahwa kedua pihak
menyepakati sebuah kunci yang sama, tetapi tidak menjamin \emph{dengan siapa}
kunci itu disepakati. Dua hal itu berbeda, dan perbedaannya justru menjadi
seluruh isi serangan ini.

Karena itu, penangkalnya bukan mengganti bilangan prima atau memperpanjang
modulus --- memperbesar $p$ tidak menghalangi Mallory sedikit pun, sebab ia tidak
pernah perlu memecahkan logaritma diskret. Yang diperlukan adalah
\textbf{otentikasi}: setiap pihak harus dapat memastikan bahwa nilai $g^a$ atau
$g^b$ yang diterimanya benar-benar berasal dari orang yang dimaksud. Beberapa
cara yang dipakai dalam praktik:

\begin{itemize}
    \item \textbf{Tanda tangan digital atas nilai pertukaran.} Alice mengirim
          $g^a$ beserta tanda tangannya, dan Bob memverifikasi tanda tangan itu
          dengan kunci publik Alice yang sudah dipercaya. Mallory tidak dapat
          memalsukannya tanpa kunci privat Alice. Pembahasan tanda tangan
          digital ada pada Bab~\ref{chap:tanda-tangan-pki}.
    \item \textbf{Sertifikat dan infrastruktur kunci publik.} Sebelum pertukaran
          dimulai, kedua pihak saling menyerahkan sertifikat yang mengikat
          kunci publik kepada identitas. Inilah yang dilakukan pada TLS: kunci
          publik peladen disahkan oleh otoritas sertifikat, sehingga klien dapat
          memastikan bahwa dirinya berbicara dengan peladen yang benar.
          Pembahasannya juga ada pada Bab~\ref{chap:tanda-tangan-pki}.
    \item \textbf{Protokol Station-to-Station.} Protokol yang dirancang Diffie,
          van Oorschot, dan Wiener ini memadukan pertukaran kunci dengan tanda
          tangan atas nilai pertukaran yang telah diikatkan pada sesi tertentu,
          sehingga setiap upaya penggantian nilai dapat dideteksi
          \citep{diffie1992}.
    \item \textbf{Saluran yang telah diautentikasi sebelumnya.} Bila kedua pihak
          telah berbagi kunci jangka panjang --- misalnya melalui pertemuan
          tatap muka sekali saja --- kunci itu dapat dipakai untuk mengautentikasi
          pertukaran Diffie--Hellman berikutnya. Inilah pola yang dipakai pada
          SSH dan pada protokol Signal.
\end{itemize}

Pelajaran umumnya penting untuk diingat: pertukaran kunci yang bersifat
anonim tidak pernah cukup, sekalipun secara matematis sempurna. Kerahasiaan
memerlukan kepastian identitas, dan keduanya harus dipenuhi secara terpisah.

\subsection{Serangan Khusus terhadap Diffie-Hellman}
\label{subsec:serangan-khusus-dh}

Selain MITM, terdapat sejumlah serangan yang menyerang pemilihan parameter atau
penerapan protokolnya. Serangan-serangan ini tidak merusak Diffie--Hellman
sebagai gagasan, tetapi menunjukkan bahwa penerapan yang tampak benar pun dapat
menjadi tidak aman.

\begin{table}[htbp]
    \centering
    \begin{tabular}{@{}>{\raggedright\arraybackslash}p{3.4cm}
                       >{\raggedright\arraybackslash}p{5.4cm}
                       >{\raggedright\arraybackslash}p{4.6cm}@{}}
    \toprule
    \textbf{Serangan} & \textbf{Prinsip} & \textbf{Penangkalan} \\
    \midrule
    \textit{Subgroup confinement} &
    Bila orde $g$ memiliki faktor kecil, penyerang dapat mengarahkan kunci ke
    subgrup kecil dan mencarinya dengan menyeluruh &
    Pakai safe prime $p = 2q+1$ dan pastikan $g$ membangkitkan subgrup berorde
    $q$ \\
    \addlinespace
    Kunci publik tak sah &
    Nilai $A$ atau $B$ yang diterima tidak diperiksa apakah berada di dalam grup
    yang benar &
    Tolak nilai $A \le 1$ atau $A \ge p-1$; untuk kurva eliptik, periksa apakah
    titiknya berada pada kurva \\
    \addlinespace
    Pembangkit acak yang lemah &
    Kunci privat $a$ atau $b$ berasal dari pembangkit acak yang keluarannya dapat
    diduga, sehingga dapat didaftar seluruhnya &
    Pakai pembangkit acak yang aman secara kriptografis; jangan membangkitkan
    kunci sendiri dengan cara yang tidak teruji \\
    \addlinespace
    Penurunan mutu parameter &
    Penyerang menurunkan negosiasi agar kedua pihak memakai modulus yang lemah,
    lalu memecahkannya (Logjam 2015) &
    Tolak modulus di bawah batas minimum; pakai kelompok bernama yang baku \\
    \addlinespace
    Pohlig--Hellman &
    Bila $p-1$ seluruhnya tersusun dari faktor prima kecil, logaritma diskret
    dapat dipecah menjadi banyak masalah kecil &
    Pakai safe prime sehingga $p-1 = 2q$ dengan $q$ prima besar \\
    \bottomrule
    \end{tabular}
    \caption{Lima serangan yang menyasar parameter atau penerapan Diffie--Hellman,
    bukan matematika protokolnya}
    \label{tab:serangan-dh}
\end{table}

Perhatikan benang merah pada Tabel~\ref{tab:serangan-dh}: keempat serangan
pertama sama sekali tidak memerlukan perhitungan logaritma diskret yang berat.
Penyerangnya cukup memilih modulus yang lemah, menunggu korban memilih kunci
privatnya dengan cara yang lemah, atau memaksa korban menerima parameter yang
lemah. Inilah alasan mengapa pembakuan parameter pada
Subbagian~\ref{subsec:parameter-dh} bukan sekadar kenyamanan, melainkan
keharusan keamanan.
